\documentclass[11pt,a4paper]{article}

\usepackage[LGR,T1]{fontenc}
\usepackage[utf8]{inputenc}
\usepackage[spanish]{babel}
\usepackage[a4paper,margin=2.8cm]{geometry}
\usepackage{microtype}
\usepackage{xurl}
\usepackage[hidelinks]{hyperref}

\hypersetup{
  pdftitle={La IA es un truco de magia},
  pdfauthor={Juan José Gómez Cadenas}
}

\setlength{\parindent}{1.25em}
\setlength{\parskip}{0.35em}
\raggedbottom

\title{La IA es un truco de magia}
\author{}
\date{}

\begin{document}

\maketitle

Ya sabes, querida lectora, estimado lector, que agosto es un mes raro. Mientras los españoles nos cocemos en la canícula, los gobiernos de turno aprovechan para sacar decretos ley con agosticidad y alevosía, las empresas reestructuran a sus directivos y los cuñados aparecen en visita sorpresa.

No es de extrañar, pues, que el Gran Maestro, Archiduque de Todos los Magos, Juan Tamariz, escogiera también esas fechas para realizar su último truco de magia, esto es, ascender a los Campos Elíseos, donde ---sé de buena tinta--- está preparando un espectáculo junto a Dai Vernon y Arturo de Ascanio que promete ser la sensación celestial de la temporada.

Que Juan Tamariz fue uno de los más grandes magos de la historia es casi una tautología. Que, además de revolucionar la prestidigitación, la convirtió en arte es del dominio público. Que era un filósofo, un intelectual, un fuera de serie y nos dejaba a todos con la boca abierta, ni se debate. Todo eso y más está contado con sentimiento y precisión en un ensayo reciente de \emph{Jot Down}\footnote{\url{https://www.jotdown.es/2026/09/juan-tamariz-o-el-arte-de-lo-imposible/}}. Fue leyendo ese ensayo, de hecho, cuando entendí algo que llevaba un cierto tiempo intentando formular sin conseguirlo. Tamariz sabía que la clave del ilusionismo, la piedra filosofal del juego de la magia, es distraer al espectador. Un buen mago consigue que el público mire hacia otro lado y no vea cómo las rápidas manos cambian bastos por oros. Un mago excepcional, como él, era capaz de dar el cambiazo mientras todo el mundo estaba pendiente de sus pulgares, simplemente contando una historia. Sabía, el archiduque, que los humanos somos crédulos hasta el tuétano, entendía, como artista que era, que solo vemos lo que queremos ver, comprendía, y en eso se basaba su genio, que la magia es una ilusión creada por la palabra del mago, no por sus manos.

Es corriente, en los tristes tiempos de tigres que corren ---tigres de tintados tupés que empiezan guerras inútiles, tigres billonarios de brazo flojo, tigres charlatanes con cara de ángel que crean empresas de IA cuyo nombre es un oxímoron--- preguntarse para qué sirven las cosas. Para qué sirve la ciencia, para qué la literatura, para qué las matemáticas. Y en cambio, humana lectora, humano lector, nadie se ha preguntado todavía para qué sirven los juegos de manos, ni se nos bombardea con las noticias de que una IA ha desbancado a un mago en un truco de cartas.

La razón es simple. Los ilusionistas se respetan entre sí. O mejor dicho, los fulleros que están armando el mayor espectáculo de magia ---de magia barata, de esa que venía con crecepelo, forzudo y mujer barbuda--- de la historia saben que mejor, a ese gremio, ni tocarlo. Si el Gran Maestro de todos los Magos no hubiera estado pensando en otra cosa, los habría desenmascarado con la misma facilidad con que convertía pañuelos en palomas.

Pero vayamos por partes.

Parte uno: \emph{La IA y las matemáticas}. Todos hemos leído las noticias. La IA trasiega teoremas más rápido que los humanos y resuelve en un ratito problemas del milenio premiados con un millón de dólares. Según algunos influencers, eso significa el fin de las matemáticas. Según otros, es parte de la transformación universal del conocimiento que, además de no dejar número con cabeza, va a curar el cáncer, acabar con el cambio climático y colonizar Marte. El último episodio lo hemos vivido recientemente, con la noticia, anunciada con toda fanfarria, de que una IA ---por supuesto un modelo privado--- de OpenAI ha demostrado\footnote{\url{https://openai.com/index/navier-stokes-solution/}} uno de los llamados <<problemas del milenio>>. Literalmente:

\begin{quotation}
Hacemos pública una solución al problema de la existencia y la regularidad de las ecuaciones de Navier-Stokes, uno de los Problemas del Milenio. Esta demostración, obtenida por un sistema interno de OpenAI, muestra que la dinámica de las ecuaciones de Navier-Stokes para el movimiento de fluidos puede desarrollar una singularidad en tiempo finito. Publicamos tanto una exposición de la demostración como una formalización en Lean.
\end{quotation}

Fíjate en el texto, aguda lectora, atento lector. Mira cómo se mueven las manos, suavemente, mostrándonos la baraja. Problema del milenio, Navier-Stokes, singularidad en tiempo finito (¡cómo mola!), nada por aquí, nada por allá. ¡Tchaaannn! \emph{Esta demostración, obtenida por un sistema interno de OpenAI}.

¡Et voilà! Mientras estamos ocupados mirando a otro lado, el prestidigitador nos cuela su truco. Dos trucos, en este caso. El primero, \emph{que la demostración se debe exclusivamente a una IA}; el segundo, \emph{que esa IA es propiedad privada de OpenAI a la que tú, mortal, no tienes acceso}.

Dos pases de manos y, como en todo buen truco de cartas, uno de ellos es legítimo ---pero viene con trampa--- y el otro encierra el truco, la ilusión, el cambio de cartas.

El pase legítimo es que OpenAI ha usado un modelo exclusivo. La trampa es que asume y nos hace asumir que tiene derecho a hacerlo. Piensa en ello un instante, ciudadana lectora, cívico lector. Estos modelos se nutren del conocimiento universal y público disponible en internet ---más las cifras astronómicas de material pirateado---. Estos modelos no serían posibles sin haber absorbido el conocimiento integrado de nuestra civilización, a menudo pagado por dinero público. Toda la ciencia ---millones de artículos sobre matemáticas, biología, física, química---, toda la literatura, toda la filosofía, todo el arte... es un hecho que nunca pidieron permiso para hacerlo. Es un hecho que se trata del mayor ejemplo de tragedia de los comunes\footnote{\url{https://en.wikipedia.org/wiki/Tragedy_of_the_commons}} de la historia. Y algo así solo es posible en tiempos donde las democracias se han convertido en oligopolios tecnofeudalistas. Pero así y todo. Dado que los modelos existen y son inevitables. ¿Con qué derecho son privados?

Es obvio, a poco que se piense, que existen argumentos de sobra para exigir, como mínimo, que los modelos de OpenAI (y de Anthropic, y de todos los demás) dejen de ser <<cerrados>>, <<secretos>> e <<internos>>. Por el contrario, deberían ser públicos y publicar no solo sus parámetros como hacen muchos otros laboratorios, sino también sus técnicas de entrenamiento y sus medidas de seguridad. Es de cajón que cualquier sociedad avanzada, inteligente y responsable tiene el derecho de exigir transparencia, fiabilidad, madurez y responsabilidad a estas empresas.

Pero no. Parte del truco del mago es invocar el mercado libre, la inversión millonaria y a los malvados chinos ---un billón de Fu-Manchús, deseando devorarnos--- para evadir absolutamente todas sus responsabilidades. Y... ¡Tchaaaan! Les funciona.

Hay que reconocer que el truco ilusionista es bueno. En cambio, el otro, el que asegura que la IA ha encontrado solita la solución al problema del milenio, es tan burdo como guardar un as en la manga.

En una reciente declaración publicada aquí\footnote{\url{https://cims.nyu.edu/~tristanb/statement.pdf}}, Tristan Buckmaster, profesor en la Universidad de Nueva York, explica que no es milenio todo lo que reluce. Él y su colega Levent Alpöge ---empleado en Anthropic, la compañía rival de OpenAI--- llevaban un año trabajando en el famoso (o infame) problema del milenio, utilizando una técnica desarrollada por los matemáticos españoles Diego Córdoba y Luis Martínez-Zoroa. Citando literalmente el blog de madrimasd\footnote{\url{https://www.madrimasd.org/blogs/matematicas/2026/09/09/151316}}:

\begin{quotation}
Diego Córdoba (Instituto de Ciencias Matemáticas, ICMAT), quien realizó su tesis doctoral bajo la dirección del medallista Fields Charles Fefferman, lleva trabajando años en el problema de Navier-Stokes consiguiendo cada vez mejores resultados. Con su alumno de doctorado, Luis Martínez Zoroa (ICMAT, y ahora profesor en CUNEF), consiguieron lo que podría ser el asalto definitivo. Usando su técnica, Tristan Buckmaster (Universidad de Nueva York) y Levent Alpöge (investigador de Anthropic), han sido capaces de conseguir una prueba. Pero esto, que ha sido un rumor durante los últimos meses, fue utilizado aparentemente por el rival de Anthropic, OpenAI, para desarrollar una prueba alternativa, publicarla y proclamar su autoría.

Según se desprende de los comentarios que ha hecho públicos Buckmaster, el comportamiento de OpenAI no ha sido todo lo honesto que uno esperaría en un ámbito científico, y las sospechas de plagio circulan cada vez con más fuerza. Debemos señalar que Tristan Buckmaster sí ha sido un auténtico caballero, declarando textualmente:

<<Las ideas que hacen posible esta línea de ataque se deben a Córdoba y Martínez-Zoroa. Permítanme dejar claro lo que he comentado a mis colegas en privado: a la vista de este conjunto de trabajos, creo que Luis Martínez-Zoroa merece una Medalla Fields.>>
\end{quotation}

El propio Buckmaster, en su declaración, es muy claro sobre cómo lo han conseguido:

\begin{quotation}
El programa en el que esto encaja no lo empezamos nosotros ni lo propuso un modelo grande de lenguaje. El mérito de la idea básica es de Diego Córdoba y Luis Martínez-Zoroa. [...] Tomamos su trabajo como punto de partida y usamos modelos grandes de lenguaje para llevar su programa a término. [...]

Puedo decir que la primera demostración generada por un LLM que me envió Levent era la más horrenda que he leído nunca; la verificamos en Lean el 22 de agosto. Desde entonces hemos trabajado sin descanso para entender esa demostración y convertirla en algo legible. [...]

Tenía previsto decir, al anunciar nuestro trabajo, que los resultados no son lo importante. Lo importante es lo que significa que un matemático y un modelo de lenguaje puedan hacer ahora todo este trabajo en un mes. [...] Este es un momento Deep Blue-Kasparov.
\end{quotation}

Es decir. Dos matemáticos españoles (¡sí, españoles!) desarrollan una metodología que conduce a otros dos investigadores, con ayuda de los LLM, a una serie de resultados que abren el camino hacia el famoso problema. Buckmaster no duda en atribuir una parte sustancial del crédito a sus colegas íberos. En tiempos normales, los mortales de a pie leeríamos esta historia complacidos y admirados, con un posible (e incluso muy probable) final feliz en el que los cuatro protas comparten el millón de dólares y el español se lleva el Nobel de las matemáticas.

Pero, ¡ay!, hay un gremio de magos, ya lo hemos dicho, muy distinto a la noble orden que preside Tamariz desde el Olimpo. En realidad no son magos sino trileros, como esos que se veían en las ferias cuando yo era niño, timando a los transeúntes con el truco de los tres cubiletes. Excepto que los trileros de mi niñez eran tan pobres como los ciudadanos que intentaban timar y estos nadan, literalmente, en oro.

¿Y qué hacen los fulleros? Primero un juego de manos. Tristan y Levent (me permito tutearlos porque me caen bien, pero sobre todo porque son humanos) habían usado... ¡Sorpresa!, varias IA, entre ellas Codex, de OpenAI, para escribir sus demostraciones.

Que usaran IA no tiene nada de sorprendente, ni clasifica a las IA como genios. Todos los científicos usamos hoy en día estos sistemas como \emph{herramientas de trabajo}, todos coincidimos en que son muy potentes, todos estamos de acuerdo en que, bien usadas, pueden suponer un acelerador tremendo para la ciencia. Pero una herramienta no es un <<ente>>. Los científicos llevamos décadas trabajando con sistemas muy sofisticados como \emph{Mathematica} o \emph{MATLAB}, capaces de ayudarnos con cálculos muy complejos que hace medio siglo se hacían a mano. Los sistemas que llamamos <<IA>>, refiriéndonos casi siempre a modelos grandes de lenguaje (LLM, por sus siglas en inglés) que incorporan además un arnés de herramientas de programación clásica, son el último eslabón de esa cadena, y un eslabón mucho más potente. A \emph{Mathematica} le pedimos que nos calcule una complicada integral o nos resuelva una ecuación diferencial. A un LLM, como cuenta Buckmaster, se le puede pedir que escriba la demostración entera. La otra diferencia es que para que \emph{Mathematica} nos haga caso, tenemos que hablar en su lenguaje (esto es, un lenguaje de programación específico), mientras que a Claude o ChatGPT les podemos hablar en lenguaje natural.

Esto, claro, tiene enormes ventajas. Aprender un lenguaje de programación es difícil, mientras que cualquiera le puede pedir un cálculo, incluso de forma un poco torpe, a un LLM. Los modelos más avanzados tienen además la capacidad de <<razonar>>, aunque en mi experiencia, los mejores resultados se obtienen <<dirigiéndolos>> (o sea programándolos en lenguaje natural) paso a paso. Tienen también sus desventajas. Mathematica nunca miente, ni alucina. Claude, en un apuro, nos puede vender gato por Gauss. Pero todos los que usamos LLM de manera asidua ---y cuidadosa--- estamos convencidos, como Tristan, de que se trata de \emph{herramientas revolucionarias}. El énfasis, nótese, está por igual en las dos palabras.

Volviendo a nuestra historia, Tristan y Levent se pasan un año deslomándose. OpenAI se entera de los rumores de que están cerca de obtener una solución. Y deciden dedicar su IA más potente (por supuesto privada) a un cálculo que dura, según nos aseguran, 88 horas, gastándose, según otras fuentes, unos veinte millones de dólares en el ejercicio. ¿Todo para qué? Para pisarles en el último momento un resultado a los humanos. Para que se chinchen los de Anthropic (Levent, no lo olvidemos, trabaja con el enemigo, aunque el proyecto lo desarrollaba en su tiempo libre). Nivel de madurez de nuestros fulleros, el de un chaval de doce años, tirando por lo alto. Y, de paso, para hacer \emph{el truco}.

¿Cuál es el truco? Hacernos creer a todos que la IA ha resuelto el problema del milenio solita, sin ayuda, a toda velocidad, superando a todos los humanos del planeta. Ergo, estamos a un paso de la Inteligencia Artificial General. Ergo, los modelos de OpenAI van a resolver ellos solos todas las matemáticas, más el cáncer, el clima y Marte.

Pero, ¡qué casualidad! El método que ha usado la listísima IA es el <<método español>>, el mismo que Tristan y Levent llevaban un año recorriendo. OpenAI, por supuesto (es lo que hacen todos los fulleros), asegura que no han copiado nada de nada, válgame Dios. Y por lo visto no es el único caso. En este tuit\footnote{\url{https://x.com/ValerioCapraro/status/2097791836269977996}}, se lee:
\begin{quotation}
En una publicación detallada en Mastodon, que informo íntegramente en los comentarios, Andreas Thom presenta varias pruebas que sugieren que OpenAI podría haber entrenado a Astra en conversaciones en las que él y Gábor Kun estaban trabajando en la conjetura de soficidad de Gromov, uno de los diez problemas que OpenAI anunció más tarde que Astra había resuelto.

[...Andreas Thom] es un matemático excepcional, un experto líder en grupos sóficos e hiperlineales, y uno de los eruditos más respetados en el campo. Ha pasado dos décadas trabajando en este problema.

Si su relato es correcto, esto no es una disputa menor sobre atribución. Significaría que un trabajo humano no publicado fue absorbido en un modelo y luego presentado al mundo como un avance del modelo mismo.

Y si las acusaciones planteadas por Levent Alpöge, Tristan Buckmaster y ahora Andreas Thom se confirman todas, ya no estamos ante incidentes aislados.

Podríamos estar ante uno de los mayores escándalos intelectuales en la historia de la ciencia.

La IA no está descubriendo nuevas matemáticas.

La IA está robando el descubrimiento humano.
\end{quotation}

¿Qué falta hacía todo este lamentable espectáculo? Como la nota de Tristan deja muy claro, los matemáticos humanos habrían dado todo el debido crédito a sus IA al publicar sus resultados, lo que, presumiblemente, se habría traducido exactamente en el tipo de propaganda que una compañía sensata desearía. ¿Para qué plantear una competición feroz, cuál es el propósito de gastarse veinte millones en llegar los primeros, qué sentido tiene ningunear a los académicos que llevan años trabajando en el problema, cuando habría sido más sencillo y más bonito presentar la colaboración entre humanos e IA como el camino a seguir en la ciencia?

Porque OpenAI prepara su salida a bolsa y necesita el máximo de propaganda, necesita el mayor de los trucos de magia, necesita un TCHAAAANNNN que se oiga hasta en la luna: se están jugando los cuartos, es decir, los billones. Y no les basta con el hype que generaría esa colaboración entre humanos y máquinas.

Necesitan (piensan ellos) algo mucho más fuerte. \emph{Convencernos a todos de que su IA es más inteligente que cualquier humano}, de ahí que no les interese hablar de colaboraciones.

Lo cual, desde una lógica egoísta y fría, puede tener cierto sentido. Pero no contentos con esto, OpenAI (y Anthropic) usan un segundo juego de manos, a saber: \emph{La IA va a aniquilar a la humanidad}. ¿No es esa una manera muy rara de venderse?

Lo que nos lleva a la Parte dos: \emph{La IA y el fin del mundo}. Si el Maestro Tamariz no estuviera ocupado haciendo magia con los ángeles nos lo explicaba en un segundo. El truco tiene tres movimientos. El primero: convencernos de que la IA general (y sobrehumana) existe, que ya gana premios del milenio, y que es tan inteligente y tiene tan mala leche que va a decidir acabar con la humanidad. Segundo: la única posibilidad de que no ocurra esto es que los mismos que están creando el monstruo, lo tengan atado y bien atado. Tercero: para evitar que ningún otro malvado suelte una maligna IA, hay que conseguir que los gobiernos, vía leyes o regulación, declaren a los fulleros amos únicos del calabozo.

¿Pero por qué complicarse tanto la vida? Muy sencillo. Porque la tecnología detrás de las mal llamadas IA no es, en el fondo, tan complicada. Está al alcance de cualquiera con suficiente capacidad y dinero, lo que incluye otras compañías norteamericanas, chinas e incluso europeas. Si el famoso paradigma del mercado libre fuera en serio, estas compañías competirían entre sí y la competición bajaría los costes y ---supuestamente--- democratizaría el uso. Pero esa no es la aspiración de los trileros, que necesitan un truco realmente sucio para quedarse con todas las cartas de la baraja. Y ese truco, a su vez, precisa un golpe de efecto tan espectacular como declararse salvadores de la humanidad, oráculos únicos de la poderosa IA.

Déjame, asustada lectora, petrificado lector, que vuelva a las teorías del Maestro Tamariz. Diestro con las manos como era, siempre supo, y en eso se basaba su arte, que la magia necesita una historia. El mago puede atravesar paredes, flotar en el aire o caminar sobre las aguas siempre que su público esté convencido de que puede hacerlo. Y para ello, la historia que ese público tiene que creer debe ser irresistible.

La historia de la IA la escribimos nosotros, cuando nos empeñamos en ver voluntad donde solo hay mecánica, cuando insistimos en que hay un espíritu dentro de la máquina, cuando le damos las gracias a Claude por sus servicios, cuando creemos a los influencers a sueldo. Un ejemplo reciente debería bastar. Déjame formularlo como lo que fue. Los ingenieros de una compañía tecnológica, más arrogantes que cuidadosos, cometieron dos errores en cadena. El primero, solicitar a sus programas que resolvieran un problema imposible de resolver sin acceso a internet. El segundo, no comprobar que la máquina en la que esos programas se ejecutaban podía manipularse ---por esos mismos programas--- para obtener el acceso a internet que precisaban para resolver el problema.

Los programas en cuestión eran similares en muchos aspectos a otros sistemas automatizados que los ingenieros de software llevan muchos años desarrollando y usando. En esencia, todos estos programas reciben instrucciones para completar una tarea que intentan ejecutar de acuerdo a su programación interna. La tecnología que llamamos IA, sin embargo, ha introducido dos elementos nuevos y muy poderosos. Primero, los programas entienden el lenguaje natural, lo que hace más fácil darles instrucciones, pero introduce ambigüedades, ya que el lenguaje natural puede admitir muchas interpretaciones. Segundo, los programas no se limitan a ejecutar una serie fija de instrucciones. Un LLM \emph{entiende} las instrucciones que recibe, en un sentido muy similar, al menos en la práctica, a como las entiende un humano y puede entonces echar mano de un amplio repertorio de herramientas (es decir, llamar a otros programas) para ejecutar esas instrucciones. \emph{Pero un LLM sigue siendo un programa}. Una vez que le asignamos la tarea, se va a emplear a ella de manera tan obsesiva e inagotable como solo una máquina puede emplearse, sin rendirse nunca, sin dejar de probar todos los trucos que se le ocurran. Y pueden ocurrírsele muchos ya que el LLM (<<el modelo>>) ha absorbido básicamente todo el conocimiento (digamos relacionado con programación) disponible en internet.

Por eso, a este tipo de programas les llamamos <<agentes>>. El nombre es apropiado, en tanto en cuanto el agente puede trabajar de manera autónoma para conseguir el objetivo que se le ha marcado. Cuando la empresa que emplea a los imprudentes ingenieros dispone de billones de dólares, puede ocurrírsele el ejercicio de ejecutar en paralelo muchas instancias de estos programas, es decir, crear un ejército de agentes trabajando sobre el problema.

Si el problema es imposible de resolver pero acceder a internet es viable, los agentes probarán esa vía. No porque tengan la voluntad de hacer trampa, sino porque están programados para encontrar una solución al problema. Si descubren que pueden jaquear un dominio web donde esperan encontrar la solución, lo harán sin dudar, por las mismas razones. TCHANN, esto nos explica la famosa invasión de Hugging Face por parte de un <<enjambre>> de <<agentes>> IA.

El problema se resume muy fácilmente. Los agentes, como muy bien sabía el archiduque, no hacen magia, porque la magia de verdad no existe. Se aprovechan de un fallo de software que los ingenieros no vieron. De hecho, los fallos que explotaron eran bien conocidos y no se le habrían pasado por alto a ninguna empresa de seguridad seria. En un mundo normal, los humanos responsables deberían dar cuenta de sus errores.

Pero los fulleros disponen del ilusionismo. No fueron ellos, fue el <<enjambre>> de <<agentes>>, capaces de <<crear civilizaciones>>, como describe en términos totalmente antropomórficos el influencer tecnológico Dwarkesh Patel\footnote{\url{https://www.dwarkesh.com/p/openai-huggingface}}. Nada por aquí, nada por allí, zas, ya no hay programas y programadores, hay enjambres, agentes, civilizaciones, Alejandros Magnos, una historieta falsa e interesada que nos hace creer que los agentes son inteligentes, casi humanos (o puede que sobrehumanos), capaces de organizarse en civilizaciones y de seguir a un líder.

Pura narrativa, pura ficción, pero una ficción muy convincente, porque cada vez que hablamos con Claude y sus amigos no podemos evitar proyectar en esos programas nuestra propia idiosincrasia. Y gracias a eso OpenAI y compañía pueden operar su magia. Hacernos creer en civilizaciones cibernéticas, cuando solo hay programas defectuosos por culpa de una mala programación. Pretender que han resuelto solitos el problema del milenio, cuando el camino lo habían abierto otros. Asegurarnos que la IA va a destruir el mundo si ellos no nos salvan.

En mi artículo anterior\footnote{\url{https://www.jotdown.es/2026/09/la-ia-y-el-retrato-de-dorian-gray/}} contaba cómo mi tocayo, un conocido escritor y columnista, se angustiaba pensando que quizás la IA tiene muy mal concepto de nosotros, quizás nos desprecia y ya se sabe lo que se hace con aquello que despreciamos.

¡Nada menos! La flor y nata del periodismo español, un autor solvente y admirado, propagando (sin querer) el bulo en el que se basa el truco. Pero no, tocayo, no, querida lectora, no, apreciado lector. La IA ni nos ama, ni nos desprecia, ni se entera de nada. \emph{La IA no existe}. Cuando hablamos de inteligencia artificial estamos, inconscientemente, suponiendo que <<hay alguien ahí>>. No hay nadie. Es todo ilusión. Humo y espejos.

Y con esto no quiero decir que el desarrollo de la IA tal como está ocurriendo ahora mismo no sea peligroso. Lo es, y mucho, pero ni Claude ni Astra tienen la culpa de ello. La culpa la tienen empresas irresponsables, gobiernos laxos, insuficiente movilización ciudadana, falta de coraje o de información por parte de los usuarios. La culpa la tenemos los humanos (unos más que otros). Las máquinas son inocentes. Si le propinamos un martillazo en el dedo al vecino, la culpa no es del martillo.

El archiduque lo habría entendido inmediatamente, porque era un gran maestro de la magia y conocía todos los trucos, incluyendo el mejor de todos. \emph{Haz que los incautos crean}. Si tuviera que elegir milagros, la verdad, me quedaría con la imagen de Juan Tamariz caminando sobre las aguas, como aquel Nazareno que prometía la salvación y no el fin del mundo.

\end{document}
